\documentclass{article}

% Encoding
\usepackage[utf8]{inputenc}
\usepackage[T1]{fontenc}

% AMS
\usepackage{amsthm}
\usepackage{amsmath}
\usepackage{amsfonts}
\usepackage{amssymb}

% Graphics and tables
\usepackage{graphicx} 
\usepackage{graphbox}
\usepackage{float}
\usepackage{booktabs}
\usepackage{multirow}
\usepackage{rotating}
\usepackage{array}
\usepackage{xcolor}
\usepackage{subcaption}
\usepackage[ruled]{algorithm2e}

\usepackage{booktabs}
\usepackage{multirow}

% Comments
\usepackage{comment}

% New column types
\newcolumntype{C}[1]{>{\centering\arraybackslash}p{#1}}
\newcolumntype{R}[1]{>{\raggedleft\let\newline\\\arraybackslash\hspace{0pt}}m{#1}}
\newcolumntype{L}[1]{>{\raggedright\let\newline\\\arraybackslash\hspace{0pt}}m{#1}}

% Break citations
\usepackage{breakcites}

% No widows or orphans (use \nowidow[3] or \noclub[3] at the end of paragraph, respectively)
\usepackage{nowidow}

% For reference enumerations lists
\usepackage{enumitem}

% Margins setup
\usepackage[top=1.5cm,bottom=2cm,right=2.5cm,left=2.5cm]{geometry}
\usepackage{titling}

% References
\usepackage{natbib}

% Check the OS
\usepackage{ifplatform}

% Symbols
\usepackage{textcomp}
\usepackage{bbm}

% pdf version 1.7
\ifwindows
	\pdfoptionpdfminorversion 7
\fi

% No indentation in the beginning of the paragraph
% \setlength\parindent{0cm} 

% Allow math display page break
\allowdisplaybreaks

% Macros
\newcommand{\lp}{\left(}
\newcommand{\rp}{\right)}
\newcommand{\lc}{\left[}
\newcommand{\rc}{\right]}
\newcommand{\lb}{\left\{}
\newcommand{\rb}{\right\}}
\newcommand{\lf}{\left.}
\newcommand{\ri}{\right.}
%
\newcommand{\R}{\mathbb{R}}
\newcommand{\T}{\mathbb{T}}
\newcommand{\Sp}{\mathbb{S}}
\newcommand{\Z}{\mathbb{Z}}
\newcommand{\N}{\mathbb{N}}
%
\newcommand{\rd}{\mathrm{d}}
\newcommand{\cL}{\mathcal{L}}
\newcommand{\cA}{\mathcal{A}}
\newcommand{\bpsi}{\boldsymbol{\Psi}}
\newcommand{\buno}{\boldsymbol{1}}
\newcommand{\bl}{\boldsymbol{l}}
\newcommand{\bx}{\boldsymbol{x}}
\newcommand{\bz}{\boldsymbol{z}}
\newcommand{\bt}{\boldsymbol{t}}
\newcommand{\bs}{\boldsymbol{s}}
\newcommand{\be}{\boldsymbol{e}}
\newcommand{\bm}{\boldsymbol{m}}
\newcommand{\bn}{\boldsymbol{n}}
\newcommand{\by}{\boldsymbol{y}}
\newcommand{\bX}{\boldsymbol{X}}
\newcommand{\bY}{\boldsymbol{Y}}
\newcommand{\bZ}{\boldsymbol{Z}}
\newcommand{\bJ}{\boldsymbol{J}}
\newcommand{\bU}{\boldsymbol{U}}
\newcommand{\bR}{\boldsymbol{R}}
\newcommand{\bF}{\boldsymbol{F}}
\newcommand{\bmu}{\boldsymbol\mu}
\newcommand{\bu}{\boldsymbol{u}}
\newcommand{\bv}{\boldsymbol{v}}
\newcommand{\bb}{\boldsymbol{b}}
\newcommand{\bba}{\boldsymbol{a}}
\newcommand{\br}{\boldsymbol{r}}
\newcommand{\bh}{{\boldsymbol{h}}}
\newcommand{\bd}{\boldsymbol{d}}
\newcommand{\bk}{\boldsymbol{k}}
\newcommand{\bp}{\boldsymbol{p}}
\newcommand{\one}{\mathbf{1}}
\newcommand{\zero}{\mathbf{0}}
\newcommand{\two}{\mathbf{2}}
\newcommand{\bxi}{\boldsymbol\xi}
\newcommand{\bphi}{\boldsymbol\varphi}
\newcommand{\ba}{\boldsymbol\alpha}
\newcommand{\bbe}{\boldsymbol\beta}
\newcommand{\bbeta}{\boldsymbol\eta}
\newcommand{\bzeta}{\boldsymbol\zeta}
\newcommand{\bga}{\boldsymbol\gamma}
\newcommand{\bthe}{\boldsymbol\theta}
\newcommand{\btheta}{\boldsymbol\theta}
\newcommand{\bkappa}{\boldsymbol\kappa}
\newcommand{\bTheta}{\boldsymbol\Theta}
\newcommand{\Ical}{\mathcal{I}}
\newcommand{\bsigma}{\boldsymbol\sigma}
\newcommand{\bSigma}{\boldsymbol\Sigma}
\newcommand{\bLambda}{\boldsymbol\Lambda}
\newcommand{\bGamma}{\boldsymbol\Gamma}
\newcommand{\bDelta}{\boldsymbol\Delta}
\newcommand{\bO}{\mathbf{0}_q}
\newcommand{\bOd}{\mathbf{0}_{q-1}}
\newcommand{\Hcal}{\mathcal{H}}
\newcommand{\bHcal}{\boldsymbol{\mathcal{H}}}
\newcommand{\X}{\boldsymbol{\mathcal{X}}_{\bx,p}}
\newcommand{\W}{\boldsymbol{\mathcal{W}}_{\bx}}
\newcommand{\bnabla}{\boldsymbol\nabla}
\newcommand{\blambda}{\boldsymbol\lambda}
\newcommand{\bB}{\boldsymbol{B}}
\newcommand{\bK}{\boldsymbol{K}}
\newcommand{\bH}{\boldsymbol{H}}
\newcommand{\bA}{\boldsymbol{A}}
\newcommand{\bV}{\boldsymbol{V}}
\newcommand{\bI}{\boldsymbol{I}}
\newcommand{\bS}{\boldsymbol{S}}
\newcommand{\bP}{\boldsymbol{P}}
\newcommand{\bQ}{\boldsymbol{Q}}
\newcommand{\bT}{\boldsymbol{\mathcal{T}}}
\newcommand{\bM}{\boldsymbol{M}}
\newcommand{\bW}{\boldsymbol{W}}
\newcommand{\pr}{^{\prime}}
\newcommand{\cqfd}{\hfill $\square$}
%
\newcommand{\equald}{\stackrel{d}{=}}
\newcommand{\inlaw}{\stackrel{d}{\rightsquigarrow}}
\DeclareMathAlphabet\mathbfcal{OMS}{cmsy}{b}{n}
\DeclareMathOperator{\sech}{sech}
\DeclareMathOperator{\csch}{csch}
\DeclareMathOperator{\arcsec}{arcsec}
\DeclareMathOperator{\arccot}{arcCot}
\DeclareMathOperator{\arccsc}{arcCsc}
\DeclareMathOperator{\arccosh}{arcCosh}
\DeclareMathOperator{\arcsinh}{arcsinh}
\DeclareMathOperator{\arctanh}{arctanh}
\DeclareMathOperator{\arcsech}{arcsech}
\DeclareMathOperator{\arccsch}{arcCsch}
\DeclareMathOperator{\arccoth}{arcCoth} 
% More macros
\newcommand{\Xb}{\bX}
\newcommand{\Sb}{\mathbf{S}}
\newcommand{\thetab}{\boldsymbol\theta}
\newcommand{\etab}{\boldsymbol\eta}
\newcommand{\bnab}{\boldsymbol\nabla}
\newcommand{\xib}{\boldsymbol\xi}
\newcommand{\Vb}{\mathbf{V}}
\newcommand{\Qb}{\mathbf{Q}}
\newcommand{\Rb}{\mathbf{R}}
\newcommand{\Zb}{\mathbf{Z}}
\newcommand{\ab}{\mathbf{a}}
\newcommand{\ub}{\mathbf{u}}
\newcommand{\Tb}{\mathbf{T}}
\newcommand{\vb}{\mathbf{v}}
\newcommand{\xb}{\bx}
\newcommand{\yb}{\by}
\newcommand{\wb}{\mathbf{w}}
\newcommand{\Ab}{\mathbf{A}}
\newcommand{\Bb}{\bB}
\newcommand{\Hb}{\mathbf{H}}
\newcommand{\Ob}{\mathbf{O}}
\newcommand{\Cb}{\mathbf{C}}
\newcommand{\Db}{\mathbf{D}}
\newcommand{\Ub}{\mathbf{U}}
\newcommand{\Mb}{\mathbf{M}}
\newcommand{\Wb}{\mathbf{W}}
\newcommand{\Yb}{\mathbf{Y}}
\newcommand{\Lb}{\mathbf{L}}
\newcommand{\Gb}{\mathbf{G}}
\newcommand{\sigmad}{\sigma_{d}}
\newcommand{\sigmar}{\sigma_{\bd}}
\newcommand{\sigmai}{\sigma_{d_i}}
\newcommand{\sigmaim}{\sigma_{d_i-1}}
\newcommand{\sigmaj}{\sigma_{d_j}}
\newcommand{\sigmamj}{\sigma_{-d_j}}
\newcommand{\sigmamjk}{\sigma_{-(d_j,d_k)}}
\newcommand{\sigmak}{\sigma_{d_k}}
\newcommand{\sigmaell}{\sigma_{d_\ell}}
\newcommand{\sigmajm}{\sigma_{d_j-1}}
\newcommand{\sigmarm}{\sigma_{\bd-1}}
\newcommand{\sigmarmij}{\sigma_{d_1-1,\stackrel{i\neq j}{\ldots},d_r-1}}
\newcommand{\Sdr}{(\mathbb{S}^{d})^r}
\newcommand{\Sr}{\mathbb{S}^{\bd}}
\newcommand{\Srm}{\mathbb{S}^{\bd-1}}
\newcommand{\Sd}{\mathbb{S}^{d}}
\newcommand{\Sj}{\mathbb{S}^{d_j}}
\newcommand{\Suno}{\mathbb{S}^{d_1}}
\newcommand{\Sdos}{\mathbb{S}^{d_2}}
%\newcommand{\Smj}{\mathbb{S}^{-d_j}}
\newcommand{\Sk}{\mathbb{S}^{d_k}}
\newcommand{\Sell}{\mathbb{S}^{d_\ell}}
\newcommand{\Sjk}{\mathbb{S}^{d_j,d_k}}
\newcommand{\Smj}{\mathbb{S}^{\bd}_{-j}}
\newcommand{\Smjk}{\mathbb{S}^{\bd}_{-(j,k)}}
\newcommand{\Sjm}{\mathbb{S}^{d_j-1}}
\newcommand{\Sim}{\mathbb{S}^{d_i-1}}
\newcommand\independent{\protect\mathpalette{\protect\independenT}{\perp}}
\def\independenT#1#2{\mathrel{\rlap{$#1#2$}\mkern2mu{#1#2}}}

% Commands with arguments
\newcommand{\lrp}[1]{\left(#1\right)}
\newcommand{\lrc}[1]{\left[#1\right]}
\newcommand{\lrb}[1]{\left\{#1\right\}}
\newcommand{\lrpbig}[1]{\big(#1\big)}
\newcommand{\lrcbig}[1]{\big[#1\big]}
\newcommand{\lrbbig}[1]{\big\{#1\big\}}
\newcommand{\blrp}[1]{\big(#1\big)}
%
\newcommand{\Prob}[1]{\mathbb{P}\lb #1\rb}
%\newcommand{\E}[1]{\mathbb{E}\lp #1\rp}
\newcommand{\V}[1]{\mathbb{V}\mathrm{ar}\lc #1\rc}
\newcommand{\Es}[2]{\mathbb{E}_{#2}\lc #1\rc}
\newcommand{\Vs}[2]{\mathbb{V}\mathrm{ar}_{#2}\lc #1\rc}
\newcommand{\Cov}[2]{\mathbb{C}\mathrm{ov}\lc #1,#2\rc}
\newcommand{\supp}[1]{\mathrm{supp}\lp #1\rp}
\newcommand{\diag}[1]{\mathrm{diag}\lp #1\rp}
\newcommand{\cmod}[1]{\mathrm{cmod}\left(#1\right)}
%
\newcommand{\mse}[1]{\mathrm{MSE}\lrc{#1}}
\newcommand{\mise}[1]{\mathrm{MISE}\lrc{#1}}
\newcommand{\amise}[1]{\mathrm{AMISE}\lrc{#1}}
%
\newcommand{\df}[2]{\frac{\rd #1}{\rd #2}}
\newcommand{\dftwo}[2]{\frac{\rd^2 #1}{\rd #2^2}}
\newcommand{\pf}[2]{\frac{\partial #1}{\partial #2}}
\newcommand{\pftwo}[2]{\frac{\partial^2 #1}{\partial #2^2}}
\newcommand{\pfmix}[3]{\frac{\partial^2 #1}{\partial #2\partial #3}}
\newcommand{\norm}[1]{\left|\left| #1\right|\right|}
\newcommand{\abs}[1]{\left| #1\right|}
\newcommand{\tr}[1]{\mathrm{tr}\left[#1\right]}
\newcommand{\vect}[1]{\mathrm{vec}\left(#1\right)}
\newcommand{\vech}[1]{\mathrm{vech}\left(#1\right)}
\newcommand{\scalprod}[2]{\langle#1,#2\rangle}
%
\newcommand{\vlinel}[1]{\multicolumn{1}{|c}{#1}}
\newcommand{\vliner}[1]{\multicolumn{1}{c|}{#1}}
%
\newcommand{\alert}[1]{\textcolor{red}{#1}}

% Orders
\newcommand{\Om}[1]{\Omega_{#1}}
\newcommand{\om}[1]{\omega_{#1}}
\newcommand{\Iq}[2]{\int_{\Omega_{q}} #1\,\omega_{q}(d #2)}
\newcommand{\Iqr}[3]{\int_{\Omega_{q}\times\R} #1\,d #3\,\omega_{q}(d #2)}
\newcommand{\Iqp}[3]{\int_{\Omega_{q}\times\R^p} #1\,d #3\,\omega_{q}(d #2)}
\newcommand{\Iqm}[2]{\int_{\Omega_{q-1}} #1\,\omega_{q-1}(d #2)}
\newcommand{\Ir}[2]{\int_{\R} #1\,d #2}
\newcommand{\Iqq}[3]{\int_{\Omega_{q_1}\times\Omega_{q_2}} #1\,\omega_{q_2}(d #3)\,\omega_{q_1}(d #2)}
\newcommand{\Irr}[3]{\int_{\R\times\R} #1\,d #3\,d #2}
%
\DeclareFontFamily{OT1}{pzc}{}
\DeclareFontShape{OT1}{pzc}{m}{it}{<-> s * [1.10] pzcmi7t}{}
\DeclareMathAlphabet{\mathpzc}{OT1}{pzc}{m}{it}
\newcommand{\order}[1]{\mathpzc{o}\lp#1\rp}
\newcommand{\Order}[1]{\mathcal{O}\lp#1\rp}
\newcommand{\orderp}[1]{\mathpzc{o}_\mathbb{P}\lp#1\rp}
\newcommand{\Orderp}[1]{\mathcal{O}_\mathbb{P}\lp#1\rp}

% New counter
%\newcommand{\co}{\addtocounter{align}{1}\arabic{align}}

% New theorems
\newtheorem{definition}{Definition}[section]
\newtheorem{theorem}{Theorem}[section]
\newtheorem{corollary}{Corollary}[section]
\newtheorem{remark}{Remark}[section]
\newtheorem{proposition}{Proposition}[section]
\newtheorem{lemma}{Lemma}[section]
\newtheorem{algo}{Algorithm}[section]
\newtheorem{example}{Example}[section]
%\renewenvironment{proof}[1]{\textit{Proof#1.}}{\qed\\} 

% Roman enumerate
\renewcommand{\theenumi}{\textit{\roman{enumi}}}

% Percentage sign without typing \%
\newcommand{\pct}{\%}

% Definition
\newcommand{\defin}{:=}
\newcommand{\indef}{=:}

\title{Research diary Vinícius PhD}

\date{\today}

\begin{document}

\maketitle

\section*{2025-06-10}

TODO for next meeting:
\begin{itemize}
    \item[$\checkmark$] Define $f_0$
    \item[$\checkmark$] Define bold h
    \item[$\checkmark$] Introduce stretchers before introducing the statistic $T_n$
    \item[$\checkmark$] Use $\rho(h)$
    \item[$\checkmark$] Use independence estimator $f_0$
    \item[$\checkmark$] Do not use $i$ as a component index
    \item[$\checkmark$] Expectation of $f$ equal to expectation of $f_0$
    \item[$\checkmark$] Mention product kernel as a definition
    \item[$\checkmark$] Use \verb|\b| whenever possible
    \item[$\checkmark$] Fix $c$ (normalizing constant)
    \item[$\checkmark$] Define $L_u$ earlier (fix swapped notation)
    \item[$\checkmark$] Remove references
    \item [$\checkmark$] Remove compilation errors in the draft
    \item Aim for application of Lemma 1 in the article
    \item Use Taylor expansion (done in the lemma 1 paper)
    \item [$\checkmark$] Properly introduce the notation and the introduction of the density stretchers. Define lambda and other necessary notation as they appear in the polysphere paper.
    \item[$\checkmark$] Split sections and better organize the paper
    \item [$\checkmark$] Apply law of large numbers to $I_{n1}^{(1)}$ as we did on the board
    \item [$\checkmark$] Polysphere paper: variance of $f_0$, read sections 2–3
    \item[$\checkmark$] Add conjecture on variance of terms generalizing Lemma 7 by adding bandwidth multipliers raised to their dimensions
    \item [$\checkmark$] Write the conjecture of the main theorem, its proof using the lemmas, assuming they hold
    \item [$\checkmark$] Leave the proof of the lemma for $I_{n1}^{(2)}$ aside
\end{itemize}

\section*{2025-06-13}

TODO for next meeting:
\begin{itemize}
    \item [$\checkmark$] Re-read the paperwriting document to avoid mistakes while writing
    \item [$\checkmark$] Introduce lambda and rho at the beginning of the paper, in section 2; mention the normalizing constant
    \item [$\checkmark$]Change references to match the format used in directionalstats.bib
    \item [$\checkmark$]Add the new references to the file no-directional-stats
    \item [$\checkmark$] Be aware that I will have to introduce context in section 1 (spheres, goals of the work, etc.)
    \item [$\checkmark$] List the assumptions being made, using the polysphere paper
    \item [$\checkmark$] Mistake: I used the estimator’s variance instead of the kernel’s variance
    \item [$\checkmark$]Change notation in the estimators: $;h_j$
    \item [$\checkmark$]Add in section 2 results from the polysphere paper
    \item [$\checkmark$] Look at the 2015 hypersphere paper to see how they computed $Var(Z_{ni})$
    \item [$\checkmark$] \textbf{Put the main theorem and its proof based on calling lemmas. Conjecture the lemmas based on the 2015 paper structure, if needed. Then we prove the lemmas.}
\end{itemize}

\section*{2025-06-17}

TODO for next meeting:
\begin{itemize}
    \item \textbf{Do all the previous points from the previous meetings.}
    \item [$\checkmark$] Read the proof of the independence test in G-P et al (2015, Statistica Sinica, see the arXiv).
    \item [$\checkmark$] Update the conjectures on the lemmas based on the previous points.
    \item [$\checkmark$] Evaluate the variance of $Z_{ni}$ through the paper on CLT paper directional variables.
    \item [$\checkmark$] Finish the proof of the main theorem (show that $N_n$ is dominant).
    \item [$\checkmark$] Correct variance of the last two terms.
\end{itemize}

\section*{2025-06-25}
\begin{itemize}
    \item [$\checkmark$] Adopt the notation used in García-Portugués et al (2015) to decompose the statistic in the three terms in the independence test chapter.
    \item [$\checkmark$] Add the asymptotic equivalence $\int_{\Sr} L_{\bh}^k(\bx,\by) g(\by)\,\rd\by\asymp  \dots$ as a lemma.
    \item [$\checkmark$] Recalculate the variance of $I_{n1}^{(1)}$ (which will have a new name after notation change) using the lemma mentioned in the previous item, using as basis the calculations done by Eduardo in the notes.
    \item [$\checkmark$] Note that the results of this project must be coherent with the results present in García-Portugués et al (2015), considering that the directional-linear can be viewed as a particular case ($r=2$, $(d_1,d_2)=(q, 1)$ and $(h_1,h_2)=(h,g)$ --- in this case $\rho(\bh)=h^q g$).
\end{itemize}

\section*{2025-07-15}

\begin{itemize}
    \item [$\checkmark$] Put conjecture Lemma 3.2
    \item [$\checkmark$] Prove the hypothesis of the lemma of Hall, used in the proof of Lemma 3.1.
    \item Proving and presenting in a clear form Theorem 3.1
    \item Advance in the proofs of Lemmas 3.5 and 3.6
    \item [$\checkmark$] Rewrite the theorem using $T_n$ notation.
\end{itemize}

\end{document}